\subsection{Эволюция нейронных сетей через расширение топологий}

\subsubsection{Постановка задачи и мотивация}

В своей работе \cite{stanley2002neuroevolution} Stanley и Miikkulainen рассматривают нейроэволюцию как альтернативу классическим методам обучения с подкреплением, где вместо оценки функции ценности или градиентов политики напрямую реализуется эволюция поведения агента, заданного параметрами нейросети.
Ключевой вопрос: можно ли получить преимущество, эволюционируя топологию сети вместе с весами, по сравнению с традиционными подходами, где оптимизируются только веса при фиксированной архитектуре. Более подробно идейную и техническую составляющие данного подхода мы рассмотрим далее.


\subsubsection{Связь с задачей моделирования когнитивного агента}

Результаты Stanley и Miikkulainen \cite{stanley2002neuroevolution} обосновывают выбор NEAT как базового нейроэволюционного алгоритма для моделирования когнитивного агента, потому что NEAT естественно работает в задачах обучения с подкреплением, где поведение агента задаётся нейросетью, и позволяет избежать ручного подбора архитектуры.

Это даёт чёткую исследовательскую гипотезу:
  
"Эволюция топологии нейросети когнитивного агента методом NEAT может обеспечить более высокую эффективность обучения (скорость, устойчивость, компактность модели) по сравнению с классическим RL-подходом при использовании фиксированной архитектуры сети."
